% Author: Prof. Dr. Matthias Jung, DL9MJ
% Year: 2026
\begin{circuitikz}
    \tikzset{german open switch/.style={
        cute open switch, 
        bipoles/cuteswitch/thickness=0.2,
        bipoles/cuteswitch/shape=circ,
        bipoles/cuteswitch/width=0.75,
    }}
    \draw(0,0) -- ++(1,0);
    \draw[fill=DARCblue!50]  (1,-1) rectangle ++(2,2);
    \draw[fill=DARCred!50]   (3,-1) rectangle ++(1,2);
    \draw[fill=DARCblue!30]  (4,-1) rectangle ++(2,2);
    \draw(6,0) -- ++(1,0) coordinate(c0);
    \draw (2.0,0) node[](){\Large N};
    \draw (3.5,0) node[](){\Large P};
    \draw (5.0,0) node[](){\Large N};
    \draw (2.0,-1.25) node[](){Emitter};
    \draw (3.5,-1.25) node[](){Basis};
    \draw (5.0,-1.25) node[](){Kollektor};
    %
    \draw(0,0)
        to [short, *-] ++(0,2)
        to [battery2, invert, v<={$U_{BE}$}] ++(1.75,0)
        to [german open switch] ++(1.75,0)
        to [short] ++(0,-1);
    \draw(0,0)
        to [short, -*] ++(0,-3) node [rground](){}
        to [lamp] ++(3.5,0)
        to [battery2, invert, v<={$U_{CE}$}] ++(3.5,0)
        to [short] (c0);
\end{circuitikz}
